% Generated draft: requires mathematical and source review.
\documentclass[11pt]{article}
\usepackage[margin=25mm]{geometry}
\usepackage{fontspec}
\setmainfont{Latin Modern Roman}
\usepackage{amsmath,amssymb,mathtools,unicode-math}
\setmathfont{Latin Modern Math}
\usepackage{enumitem,needspace,xurl,hyperref}
\setlength{\parindent}{0pt}
\setlength{\parskip}{5pt}
\setlength{\emergencystretch}{3em}
\newcommand{\N}{\mathbb{N}}
\newcommand{\Z}{\mathbb{Z}}
\newcommand{\Q}{\mathbb{Q}}
\newcommand{\R}{\mathbb{R}}
\newcommand{\C}{\mathbb{C}}
\newcommand{\F}{\mathbb{F}}
\newcommand{\A}{\mathbb{A}}
\newcommand{\PP}{\mathbb{P}}
\newcommand{\E}{\mathbb{E}}
\newcommand{\eps}{\varepsilon}
\newcommand{\dd}{\,\mathrm{d}}
\begin{document}
\section*{UnsolvedMath ID: 30001766}
\textbf{Source problem code:} catalogue entry
\subsection{Definitions and mathematical statement}
The supplied record does not specify the coefficient field, explicitly declare that \(p\) is prime, or define a numerical measure of ``size'' or ``distance from projectivity.'' The following gives a conditional framework for the questions; these missing specifications remain unresolved.

Let \(p\) be a prime, let \(n\geq 2\) be an integer with \(p\mid n\), and let \(F\) be a specified field of characteristic \(p\). Write \(S_n\) for the symmetric group on \(n\) letters. Under the usual multilinear convention, define \(L=\operatorname{Lie}_p(n)\), also denoted \(\operatorname{Lie}(n)\) in the supplied record, to be the subspace spanned by Lie monomials in which each generator occurs exactly once in the free Lie algebra over \(F\) on \(x_1,\ldots,x_n\). The group \(S_n\) acts by permuting the generators. The record does not specify its coefficient-field or module-action conventions, so this convention requires confirmation against the source.

A finite-dimensional \(FS_n\)-module is projective if it is a direct summand of a finite free \(FS_n\)-module. Decompose
\[
L\cong P\oplus L^{\mathrm{pf}},
\]
where \(P\) is the sum of the projective indecomposable summands and \(L^{\mathrm{pf}}=\operatorname{Lie}_p^{\mathrm{pf}}(n)\) has no nonzero projective direct summand. The isomorphism type of \(L^{\mathrm{pf}}\) is determined by that of \(L\).

For a finite group \(G\) and an indecomposable \(FG\)-module \(M\), a vertex is a \(p\)-subgroup \(Q\leq G\), minimal under inclusion, such that \(M\) is a direct summand of \(\operatorname{Ind}_Q^G\operatorname{Res}_Q^G M\). A source at \(Q\) is an indecomposable \(FQ\)-module \(U\) occurring as a direct summand of \(\operatorname{Res}_Q^G M\) such that \(M\) is a direct summand of \(\operatorname{Ind}_Q^G U\). An \(FQ\)-module \(U\) is endo-permutation if \(\operatorname{End}_F(U)\), with action \((q\cdot f)(u)=qf(q^{-1}u)\), has an \(F\)-basis permuted by \(Q\).

When \(n=p^k\) with \(k\geq 1\), identify the \(n\) letters with \(\mathbb F_p^k\) and let
\[
E_n=\{v\mapsto v+a:a\in\mathbb F_p^k\}\leq S_n.
\]
Thus \(E_n\cong(C_p)^k\) acts regularly, meaning transitively with trivial point stabilizers. The choice of identification determines this subgroup up to conjugacy in \(S_n\).

The requested statement is a collection of classification questions. For the intended coefficient field, determine the size of \(L^{\mathrm{pf}}\) and describe how far \(L\) is from being projective. Dimension \(\dim_F L^{\mathrm{pf}}\) is one possible interpretation of size, but the record neither selects it nor specifies an invariant measuring distance from projectivity. For every prime-power degree \(n=p^k\), ask whether \(L^{\mathrm{pf}}\) is indecomposable, whether its vertex is conjugate to \(E_n\), whether its sources are endo-permutation modules, and what their isomorphism types are. The vertex and source questions for the whole projective-free part presuppose its indecomposability; if it is decomposable, those questions must instead be specified for its indecomposable summands. For \(p\mid n\) with \(n\) not a power of \(p\), determine the corresponding behavior of the projective-free part and its indecomposable summands. The record gives no more specific target in this last case. These are questions, not assertions of indecomposability, a particular vertex, or the endo-permutation property.
\subsection{Short English statement}
For the Lie module of \(S_n\) in characteristic \(p\), with \(p\mid n\), describe its nonprojective summands. In prime-power degree, investigate indecomposability, a regular elementary abelian vertex, and the existence and isomorphism types of endo-permutation sources; describe the changes in other degrees. The coefficient field and the meanings of size and distance from projectivity remain unspecified. This is a statement-only draft requiring human mathematical review.
\subsection{Sources}
\begin{itemize}
\item Local UnsolvedMath registry, permanent ID \(30001766\), title: \emph{Projective-Free Lie Modules of Symmetric Groups}; collection: Oberwolfach Reports Open Problems; exact source problem code: \texttt{catalogue entry}. The supplied record gives proposed year 2011 and names no proposer.
\item Supplied source citation: \emph{Mini-Workshop: Modular Representations of Symmetric Groups and Related Objects} (2011), Oberwolfach Reports, pp. 1177--1205, DOI \texttt{10.4171/owr/2011/21}. \url{https://doi.org/10.4171/owr/2011/21}. The background describes these reports as records of mathematical workshops and their open-problem discussions; the literature-source label is ``Modular Representations of Symmetric Groups and Related Objects, OWR 21/2011.'' No problem-specific page or named contributor is supplied.
\item Supplied literature reference: Lim--Tan--Willems, ``vertices of Lie modules.'' \url{https://arxiv.org/abs/1309.2085}.
\item Supplied literature reference: Erdmann--Lim--Tan, ``complexity of the Lie module.'' \url{https://arxiv.org/abs/1503.01545}.
\item Literature assessment dated exactly \texttt{2026-08-21}: the supplied curation labels the problem ``partially solved'' and states that results on complexity and vertices reduce substantial cases to \(p\)-power degrees, while the requested source and isomorphism classifications remain incomplete in the stated generality. This is supplied dated curation, not an independently verified assessment. The background explicitly requests independent verification of time-sensitive status and cited claims.
\item The supplied statement-status label is \texttt{corrected\_verified}; its verification note reports a check against the official Oberwolfach report and says that the clean statement preserves the hypotheses, notation, and open target. This is a supplied verification claim; no source was independently visited or checked for this draft.
\end{itemize}
\end{document}
